% ============================================================
% templates/t06-figures.tex   —— T6 两图并排
%
% 适用场景：两张需要对比的图（前后对照、方案 A/B、原图与结果），
%          等高、顶端对齐、水平居中。
% 提供    ：\DeckFigures
% 依赖    ：\includegraphics（graphicx，在 styles/ 里已加载）
% 参数    ：#1 左图路径 #2 右图路径 #3 左图高度 #4 右图宽度
% 容量    ：两图宽度之和**必须明显小于**版心宽，留 0.3cm 以上余量。
%           版心宽用 \SlideTextWidth 查（由 styles/ 末尾的实测段定义）
% 易错    ：① 只差 0.03cm 就会让第二张图掉到下一行，并报 **overfull vbox**
%             ——报的是竖直溢出，看不出是宽度问题，极难查。留足余量
%           ② 不要用 beamer 的 columns：它的栏间距会与给定宽度叠加，
%             试三次都算不准
%           ③ 等高要靠显式 width：\includegraphics[height=H] 的宽度按
%             原图比例推算。两张图比例不同又都想等高时，必须给其中一张
%             显式 width（= H × 该图宽/该图高），调用方算好传入
%           ④ 原图先缩放：手机拍的动辄 4000+ 像素，直接入稿会让 PDF
%             膨胀到几十 MB。先缩到长边约 2400px
%
% 【填满的样子】（两图宽度之和 4.5 + 7.0 = 11.5cm，明显小于版心宽）
%   \DeckFigures{figures/before.png}{figures/after.png}{4.5cm}{7.0cm}
% ============================================================

% --- 两图并排 ---
% 不用 beamer 的 columns：它的栏间距会与给定宽度叠加，很难算准。
% 改成两图放在同一行、用 \hfill 撑开，宽度完全可控。
\DeckDefine{\DeckFigures}[4]{%
  \par\noindent
  \includegraphics[height=#3]{#1}%
  \hfill
  \includegraphics[width=#4]{#2}%
  \par
}
